% Mục 2: Nền tảng thí nghiệm. Nguồn: results/apple-m5/machine.txt, results/apple-m5/campaign.log,
% include/platform.hpp, Makefile.

Đề yêu cầu số nền tảng shared-memory không ít hơn số thành viên. Báo cáo này có số đo của
\PCval{cross}{nhosts} máy, \PCname{apple-m5}. Mỗi thành viên thêm máy của mình bằng cách chạy
\texttt{make bench} trên máy rảnh, gửi lại thư mục \texttt{results/<tên máy>/} và thêm tên máy vào
danh sách máy của báo cáo. Bảng cấu hình (Bảng~\ref{tab:cross:machines}) và các hình của máy mới được
sinh lại tự động. Mọi nhận định về hiệu năng dưới đây chỉ nói về \PCname{apple-m5}.
\TODO{nhóm có mấy thành viên, và máy của các thành viên còn lại (sau khi họ chạy make bench)}

Máy \PCname{apple-m5} là một Apple M5 với 10 core không đồng nhất, gồm 4 P-core (performance) và 6
E-core (efficiency). Mỗi cụm có cache L2 riêng. Máy không có SMT. Ba đặc điểm của máy ảnh hưởng tới kết
quả. Thứ nhất, khi số thread vượt 4, một phần thread chạy trên E-core chậm hơn, và cache line dùng chung
phải đi qua lại giữa hai cụm. Thứ hai, macOS không cho gắn thread vào một core cố định, nên loại core
của mỗi thread do scheduler quyết định và có thể đổi giữa các lần chạy. Nhóm giảm biến thiên này bằng
cách lặp lại và lấy trung vị. Thứ ba, AArch64 có memory model yếu hơn x86-TSO. Các memory order khác
nhau sinh ra các lệnh máy khác nhau (store \texttt{relaxed} là \texttt{str}, store \texttt{release} là
\texttt{stlr}). Vì vậy lỗi chọn sai memory order có thể quan sát được trên máy này, trong khi trên x86 nó
thường không bao giờ lộ ra. CPU có các lệnh atomic LSE của ARMv8.1, nên \texttt{fetch\_add} và
\texttt{compare\_exchange} là một lệnh máy (\texttt{ldadd}, \texttt{cas}) chứ không phải vòng LL/SC.

Hệ điều hành báo cache line, đơn vị dữ liệu mà phần cứng giữ nhất quán giữa các core, dài 128 byte.
Mã nguồn đệm mọi biến dùng chung theo hằng \texttt{pc::kCacheLine}, bằng 128 trên Apple arm64 và 64
trên nền tảng khác. Nhóm kiểm tra rằng kích thước và căn chỉnh của mọi kiểu lock và của nút hàng đợi
MCS đều bằng 128 byte trên máy này.

Phần mềm gồm macOS 27.0.1 và Apple clang 21.0.0 với cờ \texttt{-std=c++20 -O2}. Apple clang không nhận
\texttt{-fopenmp}, và macOS không có \texttt{pthread\_spinlock\_t} hay \texttt{pthread\_barrier\_t}.
Vì vậy các baseline trên máy này là \texttt{std::mutex}, \texttt{pthread\_mutex\_t}, \texttt{os\_unfair\_lock},
\texttt{std::barrier} và barrier dựng từ mutex và condition variable.

Toàn bộ số đo đến từ một đợt đo duy nhất dài 1 giờ 26 phút, gồm 12\,152 lần chạy. Mỗi cấu hình có
7 lần lặp, và thứ tự các cấu hình xáo trộn trong từng lượt. Máy cắm sạc, chế độ tiết kiệm điện tắt, và
không có tiến trình đo hay build nào khác chạy cùng lúc. Load average một phút lúc bắt đầu là 6,73,
trong khoảng thường thấy của máy khi chỉ có ứng dụng nền. Không lần chạy nào sai kết quả kiểm tra.
Có 9 lần vượt thời hạn, đều là cấu hình có MCS khi số thread vượt số core, và chúng được đo lại với số vòng
lặp nhỏ hơn 100 lần (mục Phương pháp đo).
